\begin{table}
\caption{Accuracy comparison of LLaMA 7B/13B with two different tuning granularity of \ours. Columns \textbf{m} and \textbf{V} designate the modules with tunable magnitude and directional components, respectively. Each module is represented by its first letter as follows: (Q)uery, (K)ey, (V)alue, (O)utput, (G)ate, (U)p, (D)own.}
% \vspace{0.8em}
\vskip 0.1in
\setlength{\tabcolsep}{0.12mm}
\centering
\tiny
% \renewcommand\arraystretch{0.8}
\resizebox{0.46\textwidth}{!}{
\begin{tabular}{cccccc}
\toprule
\textbf{Model} & \textbf{PEFT Method} & \# \textbf{Params (\%)} & \textbf{m}  & \textbf{V} & \textbf{Avg.} \\ \hline
\multirow{3}{*}{LLaMA-7B}    & LoRA & 0.83 & - & - & 74.7  \\ 
&\ours~(Ours) & 0.84 & QKVUD& QKVUD & 78.1  \\ 
& \ours~(Ours) & 0.39 & QKVOGUD & QKV & 77.5  \\  \midrule

\multirow{3}{*}{LLaMA-13B}    & LoRA & 0.67 & - & - & 80.5 \\
& \ours~(Ours) & 0.68 & QKVUD & QKVUD & 81.5  \\ 
& \ours~(Ours) & 0.31 & QKVOGUD & QKV & 81.3  \\ \bottomrule
\end{tabular}}
% \vspace{-0.8em}
% \vspace{-0.5em}
\label{tab:llama_commonsense_tunning_granularity}
\vskip -0.1in
\end{table}


